\section{Methods}
\label{sec:methods}

We use independently seeded runs in five geometries: random density-$1/2$
bulk states, axial and tilted half-active bands, random active disks or
absorbing holes, small seeds in an absorbing background, and walls between
the two absorbing states. One step is one simultaneous lattice update.
The bit-parallel CPU simulator and dense CUDA backend use position- and
time-indexed Philox random numbers~\cite{salmon2011}. Their integer updates,
profiles and transferred checkpoints agree bit for bit in the archived
validation cases. The full protocols and computational checks are in
Supplemental Secs.~\ref{S-sec:protocols} and \ref{S-app:numerics}.

For the extended band measurements, a normal $(a,b)$ defines the coordinate
$s=(ay-bx)\bmod L$, with normal spacing $1/\sqrt{a^2+b^2}$.
We smooth over nine physical lattice spacings and follow the two exterior
profile crossings. Half the rate of change of their separation defines
the normal velocity. A complete declared cohort contributes to a window
only if every run retains the separated-front geometry throughout.
Interior activity is checked separately: the loss of a contour in a
dissolving band is not interpreted as coherent front retreat.

We integrate each bulk trajectory in time and weight independent runs
equally, including absorbed trajectories through their known absorbing
state. First recorded half-density loss and complete consensus are distinct
events; lifetime summaries retain censoring. Droplet growth means reaching
the finite box boundary before extinction and before the stated horizon,
not survival forever. Escapes are separately counted in seed experiments.

Independent runs, not lattice sites or output records, are the sampling
units. The large-system ensembles have four runs per point. Percentile
bootstrap intervals describe the observed run variation conditional on the
chosen estimator and window; with four runs their nominal $95\%$ level
does not ensure calibrated coverage. We display individual paired changes
for delicate drift claims and report small-sample reference calculations
separately. Geometry, fitting choices, finite size and finite time remain
systematic uncertainties rather than components of a pooled error bar.
